%% ch-meta-functoids.tex -- Part 2 (\ref{part:meta}), chapter 3: the hand-written argument for functoids.
\providecommand{\DS}[1]{\hyperref[ds:#1]{DS-#1}}

\chapter{Functoids}
\label{ch:meta-functoids}

The encoding of Chapter~\ref{ch:meta-k2} has no functoids of its own
beyond the table $\mathrm{FUN\text{-}DEF}$.  The soundness of the kernel's
functoid apparatus is therefore covered by a hand-written argument
(\DS{6}, \DS{7}).  This chapter states the argument and the hypotheses it
needs, and says how the kernel enforces each.

%% =========================================================
\section{What is eliminated}

Three kinds of head are \emph{defined heads}:
\begin{itemize}
\item a \texttt{def-functoid} head $F$ with parameters $\bar p$ and body
  $b_F$: its unfold rewrites $F(\bar p)$ to $b_F$ with no condition
  (Chapter~\ref{ch:defs});
\item a structure accessor $\mathit{ACC}$ of slot $k$: it rewrites
  $\mathit{ACC}(s)$ to $\mathrm{NTH}(k, s)$;
\item an anonymous record $R$ (a \texttt{lambdoid}) with bindings $\bar x$
  and body $b_R$, applied as \texttt{(apply-functoid $R$ $\bar t$)}.
\end{itemize}
$L_0$ is the part of the language in which none of them occurs.

\begin{dfn}[the elimination map $E$]
On a formula in which every defined head is applied to as many arguments
as it has parameters, $E$ is defined by recursion, arguments first:
\[
\begin{aligned}
E(F(t_1,\ldots,t_n)) &= E(b_F)[\bar p := E(\bar t)],\\
E(\mathit{ACC}(s)) &= \mathrm{NTH}(k, E(s)),\\
E(\texttt{apply-functoid}(R, t_1,\ldots,t_n)) &= E(b_R)[\bar x := E(\bar t)],
\end{aligned}
\]
with simultaneous capture-avoiding substitution; $E$ commutes with every
other constructor, connective and binder.
\end{dfn}

The meaning of a formula with defined heads is the meaning of its
E-image (\DS{7}).  ``$F(\bar t)$ denotes'' means ``$E(F(\bar t))$
denotes''; nothing more is asserted about $F$.

A record's domains are dropped by $E$: \texttt{functoid-beta} does not
read them, so they have no logical content.

\paragraph{Two lemmas.}  \emph{Substitution:}
$E(e[x:=t]) =_\alpha E(e)[x:=E(t)]$.  \emph{Definedness through the
body:} $F(\bar t)$ denotes exactly when $E(b_F)[\bar p := E(\bar t)]$
does.  In particular a defined head is not strict: $F(\bar t)$ may denote
while some $t_i$ does not.  This is the schematic reading of an unfolding
(\DS{34}).

%% =========================================================
\section{The proposition}

\begin{prop}[hand-written, conditional]\label{prop:meta-fun}
Assume version-X, the soundness of the functoid-free rules for $L_0$
(link~3 of Section~\ref{sec:meta-chain}, restricted to $L_0$), the
oracles and the infrastructure, and the hypotheses H1--H6 below.  For
every inference accepted by its checker, either its E-image is the
identity (one hypothesis with the E-image of the conclusion), or its
E-image is an inference between $L_0$ sequents accepted by some kernel
rule.

Consequently every theorem installed by \texttt{qed} over a grounded
graph with an empty bill has a valid E-image; since
$E(\textup{\textsc{falsity}}) = \textup{\textsc{falsity}}$, no such theorem is
\textup{\textsc{falsity}}, provided version-X is consistent.
\end{prop}

The proof goes rule by rule.
\begin{itemize}
\item \emph{An unfold} (\texttt{mac} at $F$, goal or hypothesis) and an
  accessor rewrite owe nothing; under $E$ they are the identity.  Owing
  nothing is correct because the substitution lemma holds at every
  argument, denoting or not.
\item \emph{A functor projection} $\mathit{NAME}@\mathit{ACC}$ owes the
  definedness of the instantiated list (\DS{29}); its E-image is an
  \texttt{nth-reduce} of a literal list, which owes the same (\DS{25}).
\item \emph{The view clause} certifies $\mathit{ACC}(V(r))$ only when the
  whole instantiated list is certified (\DS{28}); its E-image is then
  certified by the NTH clause.
\item \emph{The strictness walks} read a \texttt{def-functoid} application
  through its instance and an applied record through its contraction
  (\DS{21}); every other head they enter is, under $E$, a strict head of
  $L_0$.  So a term the walk reaches is at a strict position of the
  E-image.
\item \emph{The definedness certificate and the value tests} read a
  defined head through its instance (\DS{30}); by the two lemmas each
  verdict is the verdict on the E-image.
\item \emph{Every other rule} treats $F(\bar a)$ as an application with a
  constant head.  Such a head is never substituted, and the matchers match
  a constant head only by identity; so the E-image of an accepted
  inference is an accepted inference of the same rule.
\item \emph{Inside a record body} a record's variable is a placeholder for
  an arbitrary term, denoting or not.  Along the local-context walk of the
  rewriting rules, an assumption mentioning a crossed bound variable is
  dropped, and a term mentioning a record variable is not certified
  (\DS{27}).  A step inside the body is therefore valid after any term is
  substituted for the variable.
\end{itemize}

\paragraph{Guarded definitions.}  Heads defined by \texttt{def-constant},
the two recursion commands and \texttt{def-operator} are not eliminated:
their equations are axioms of $T$.  They are read as descriptions: off
every guard the head denotes nothing (\DS{31}).  This reading satisfies
every defining equation provided no two cases hold at once with
different values.  Under it such a head is strict, and the two readings
-- non-strict for \texttt{def-functoid} heads, strict for defined
functions -- are kept apart by H6.

%% =========================================================
\section{The hypotheses, and how each is enforced}

\begin{center}
\small
\begin{tabular}{@{}lp{\dimexpr0.44\linewidth-2\tabcolsep\relax}p{\dimexpr0.48\linewidth-2\tabcolsep\relax}@{}}
\toprule
& hypothesis & enforcement \\
\midrule
H1 & the parameters of a \texttt{def-functoid} are distinct symbols
  & functoid-hypotheses gate, fatal (\DS{35}) \\
H2 & a body names only heads defined before it (no recursion, no forward
  reference) & functoid-hypotheses gate, fatal \\
H3 & one definition per head; no theorem proven under an earlier
  definition survives a redefinition & functoid-hypotheses gate, fatal; the
  certificate key records the body's hash (\DS{17}) \\
H4 & the free variables of a body are among its parameters
  & functoid-hypotheses gate, fatal \\
H5 & every formula written into the graph is eliminable: a defined head
  occurs only applied, never as a term or as a quantifier instance or
  witness & refused at formula construction, by \texttt{forall-elim}, by
  \texttt{forsome-intro} and by \texttt{reflexivity} (\DS{33}) \\
H6 & no head is both a \texttt{def-functoid} and a library definition
  & functoid-hypotheses gate, fatal \\
\bottomrule
\end{tabular}
\end{center}

H4 carries weight.  With $G(x) := y$, $y$ free, the unfold rewrites
$G(0)$ to $y$ inside $\forall y$, so $\forall y.\,G(0)=y$ would be
provable; instantiated at $0$ and at $1$ it gives $0=1$.  The gate
refuses such a body.

The functoid-hypotheses gate (\texttt{functoid-gate.scm}) runs at the end
of every load, a cold load, a certified build and an extension of the
band alike.  It reads the registry of functoids and the log of
definitions, which records for each \texttt{def-functoid} the heads and
constants defined at that moment.  In a strict load a violation of any of
H1--H4 or H6 is an error, with one line naming the hypotheses violated;
under keep-going it is reported and the load continues.  The library has
no violation.

Two consequences for the library.  A body may name a
constant or a head only if it is defined earlier in the load order
(H2); an instance made by \texttt{declare-instance!} counts from the file
that makes it.  A structure species used as a class in a body, such as
$\mathrm{DIRECTED\text{-}SET}$ in $\mathrm{PTWISE\text{-}NET\text{-}LIMITS}$,
is a class constant through the class axiom of its
\texttt{declare-structure}.  H4 reads every declared structure species
as a constant, not as a free variable of the body.  A species is not a
registered constant: that would make it head-only, and a membership in
it would be refused as a formula.

%% =========================================================
\section{The J-MACN bound}

$\mathrm{J\text{-}MACN}(T)$, the local instances of the macetes of $T$
with any number of schema variables, bounds its argument list by
$\mathrm{SET}$.  The bound holds on the intended domain but is not
derivable from the body alone (\DS{47}).

It excludes no kernel step.  For an actual macete step the encoder takes
the list of encoded substitution values, each a member of
$\mathrm{EXPR\text{-}SET}$, which is a set; a literal list of sets is a
set.  The bound thus holds of every encoded kernel step, provided the
encoder produces a literal list of encoded terms, which is the
translator's obligation (K3).  $\mathrm{HYPS\text{-}IN}$ makes the same
kind of bound on a list of hypothesis sequents.
